% Options for packages loaded elsewhere
\PassOptionsToPackage{unicode}{hyperref}
\PassOptionsToPackage{hyphens}{url}
\PassOptionsToPackage{dvipsnames,svgnames,x11names}{xcolor}
%
\documentclass[
  11pt,
]{article}
\usepackage{amsmath,amssymb}
\usepackage{iftex}
\ifPDFTeX
  \usepackage[T1]{fontenc}
  \usepackage[utf8]{inputenc}
  \usepackage{textcomp} % provide euro and other symbols
\else % if luatex or xetex
  \usepackage{unicode-math} % this also loads fontspec
  \defaultfontfeatures{Scale=MatchLowercase}
  \defaultfontfeatures[\rmfamily]{Ligatures=TeX,Scale=1}
\fi
\usepackage{lmodern}
\ifPDFTeX\else
  % xetex/luatex font selection
\fi
% Use upquote if available, for straight quotes in verbatim environments
\IfFileExists{upquote.sty}{\usepackage{upquote}}{}
\IfFileExists{microtype.sty}{% use microtype if available
  \usepackage[]{microtype}
  \UseMicrotypeSet[protrusion]{basicmath} % disable protrusion for tt fonts
}{}
\makeatletter
\@ifundefined{KOMAClassName}{% if non-KOMA class
  \IfFileExists{parskip.sty}{%
    \usepackage{parskip}
  }{% else
    \setlength{\parindent}{0pt}
    \setlength{\parskip}{6pt plus 2pt minus 1pt}}
}{% if KOMA class
  \KOMAoptions{parskip=half}}
\makeatother
\usepackage{xcolor}
\usepackage[margin=2.2cm]{geometry}
\usepackage{longtable,booktabs,array}
\usepackage{calc} % for calculating minipage widths
% Correct order of tables after \paragraph or \subparagraph
\usepackage{etoolbox}
\makeatletter
\patchcmd\longtable{\par}{\if@noskipsec\mbox{}\fi\par}{}{}
\makeatother
% Allow footnotes in longtable head/foot
\IfFileExists{footnotehyper.sty}{\usepackage{footnotehyper}}{\usepackage{footnote}}
\makesavenoteenv{longtable}
\setlength{\emergencystretch}{3em} % prevent overfull lines
\providecommand{\tightlist}{%
  \setlength{\itemsep}{0pt}\setlength{\parskip}{0pt}}
\setcounter{secnumdepth}{-\maxdimen} % remove section numbering
\ifLuaTeX
  \usepackage{selnolig}  % disable illegal ligatures
\fi
\usepackage{bookmark}
\IfFileExists{xurl.sty}{\usepackage{xurl}}{} % add URL line breaks if available
\urlstyle{same}
\hypersetup{
  pdftitle={Square grid: computations and erratum},
  colorlinks=true,
  linkcolor={Maroon},
  filecolor={Maroon},
  citecolor={Blue},
  urlcolor={Blue},
  pdfcreator={LaTeX via pandoc}}

\title{Square grid: computations and erratum}
\author{}
\date{}

\begin{document}
\maketitle

\section{Exact and Lower-Bound Values for Right-Isosceles-Free Subsets
of the Square
Grid}\label{exact-and-lower-bound-values-for-right-isosceles-free-subsets-of-the-square-grid}

\textbf{Abstract.} For a positive integer n, let a(n) be the largest
size of a subset of the n×n integer grid that contains no three points
forming a right isosceles triangle. We compute a(n) exactly for n ≤ 7 by
constraint programming, confirm the known values for 8 ≤ n ≤ 11, and
report lower bounds for 12 ≤ n ≤ 16. Erratum: OEIS A271906 already
contains lower bounds for n = 12--15 (29, 32, 36, 38; R. Israel and G.
Resta, 2016). Our bounds for 12--15 therefore add nothing new, and our
values for 13--15 are weaker. Only n = 16 (value 40) is not listed
there. We do not prove optimality for n ≥ 8, and our upper bounds are
loose.

\subsection{1. Problem}\label{problem}

A \emph{right isosceles triangle} in ℤ² is a triple \{P, Q, R\} with a
right angle at P and \textbar PQ\textbar{} = \textbar PR\textbar.
Equivalently, there exist a vector u = (a, b) ≠ (0, 0) and a point P
such that the three points P, P+u, P+R(u) lie in the set, where R(a, b)
= (−b, a) is rotation by 90°.

Let S ⊆ \{0, \ldots, n−1\}² and call S \emph{admissible} if it contains
no right isosceles triangle. Define

a(n) = max\{ \textbar S\textbar{} : S admissible \}.

This is OEIS A271906, which lists a(1), \ldots, a(11) = 1, 2, 4, 6, 9,
11, 14, 17, 20, 23, 26.

\subsection{2. Method}\label{method}

\textbf{Exact model.} We introduce a Boolean variable for each grid
point and, for each right isosceles triangle T in the grid, the clause
``not all three points of T are chosen.'' We maximize the number of
chosen points with Google OR-Tools CP-SAT. The number of triangles grows
from 4 (n = 2) to 6020 (n = 11).

\textbf{Heuristic model.} We use simulated annealing over subsets. The
objective is the set size minus a penalty of λ times the number of
triangles fully contained in the current set, with λ = 1.5. The move
flips one grid point, and the temperature decreases geometrically from
1.0 to 0.02 over a fixed time budget. We record the largest admissible
set found.

\textbf{Independent verification.} Every reported set is checked by a
separate routine (\texttt{verify.py}). It enumerates all triples of
points and tests the right-isosceles condition directly by dot products
and squared lengths, without using the triangle list from the solver.
The checker rejects a planted violation and agrees with the solver on
all n ≤ 9.

\subsection{3. Results}\label{results}

\subsubsection{3.1 Exact values (proved by CP-SAT, optimality status
reported by the
solver)}\label{exact-values-proved-by-cp-sat-optimality-status-reported-by-the-solver}

\begin{longtable}[]{@{}lll@{}}
\toprule\noalign{}
n & a(n) & Status \\
\midrule\noalign{}
\endhead
\bottomrule\noalign{}
\endlastfoot
1 & 1 & optimal \\
2 & 2 & optimal \\
3 & 4 & optimal \\
4 & 6 & optimal \\
5 & 9 & optimal \\
6 & 11 & optimal \\
7 & 14 & optimal \\
\end{longtable}

These agree with OEIS A271906.

\subsubsection{3.2 Lower bounds}\label{lower-bounds}

\begin{longtable}[]{@{}
  >{\raggedright\arraybackslash}p{(\columnwidth - 8\tabcolsep) * \real{0.0508}}
  >{\raggedright\arraybackslash}p{(\columnwidth - 8\tabcolsep) * \real{0.1864}}
  >{\raggedright\arraybackslash}p{(\columnwidth - 8\tabcolsep) * \real{0.1356}}
  >{\raggedright\arraybackslash}p{(\columnwidth - 8\tabcolsep) * \real{0.3729}}
  >{\raggedright\arraybackslash}p{(\columnwidth - 8\tabcolsep) * \real{0.2542}}@{}}
\toprule\noalign{}
\begin{minipage}[b]{\linewidth}\raggedright
n
\end{minipage} & \begin{minipage}[b]{\linewidth}\raggedright
best found
\end{minipage} & \begin{minipage}[b]{\linewidth}\raggedright
Source
\end{minipage} & \begin{minipage}[b]{\linewidth}\raggedright
Upper bound (CP-SAT)
\end{minipage} & \begin{minipage}[b]{\linewidth}\raggedright
Matches OEIS?
\end{minipage} \\
\midrule\noalign{}
\endhead
\bottomrule\noalign{}
\endlastfoot
8 & 17 & CP-SAT, annealing & 19 & yes (A271906) \\
9 & 20 & CP-SAT, annealing & 28 & yes \\
10 & 23 & CP-SAT, annealing & 40 & yes \\
11 & 26 & CP-SAT, annealing & 121 & yes \\
12 & 29 & CP-SAT (1500 s), annealing (150 s) & 144 & yes: a(12)
\textgreater= 29 already in OEIS (R. Israel, 2016) \\
13 & 31 & annealing (150 s) & --- & weaker than OEIS bound a(13)
\textgreater= 32 (G. Resta, 2016) \\
14 & 33 & annealing (150 s) & --- & weaker than OEIS bound a(14)
\textgreater= 36 (G. Resta, 2016) \\
15 & 37 & annealing (150 s) & --- & weaker than OEIS bound a(15)
\textgreater= 38 (G. Resta, 2016) \\
16 & 40 & annealing (150 s) & --- & not listed in the OEIS entry we
checked \\
\end{longtable}

All sets in the table were verified by the independent checker. For n =
12, two different methods reached 29, but neither proved that 29 is
optimal.

\subsubsection{3.3 Example construction (n = 12, \textbar S\textbar{} =
29)}\label{example-construction-n-12-s-29}

Rows are shown from top (y = 11) to bottom (y = 0); \texttt{\#} marks a
chosen point.

\begin{verbatim}
. # # # # # . . # . . .
. . . . . . # . . # . .
. . . . . . . # . . # .
. . . . . . . . . . . #
# . . . . . . . . . . #
. # # . . . . . . . . #
. . . . . . . . . . . #
. . # # . . . . . . . #
. . . . . . . . . . . #
# . . # . . . . . . . #
. . # . . . . . . . . #
# # . . . . . . . . . #
\end{verbatim}

\subsubsection{3.4 Status of exact values and
density}\label{status-of-exact-values-and-density}

\begin{itemize}
\tightlist
\item
  \textbf{a(8) = 17 is now proved} (CP-SAT with 2×2 block cuts, 93 s,
  status OPTIMAL). It matches OEIS, so it is not new.
\item
  \textbf{a(12) is not proved.} After 3000 s the solver held a set of
  size 29 with upper bound 72, which is the trivial 2×2 bound. The
  solver gave no improvement on the trivial bound.
\item
  \textbf{Attempt to prove a(12) ≤ 29 with window cuts.} Adding the
  constraint that every k×k window (k ≤ 8) holds at most a(k) points,
  with the proved values a(k), did not resolve the question. For target
  30, CP-SAT returned UNKNOWN after 2400 s on 2 cores. The n = 9 target
  21 was also UNKNOWN after 200 s.
\item
  \textbf{Density.} Covering the n×n grid by k×k blocks gives limsup
  a(n)/n² ≤ a(k)/k². The best k we can compute (k = 7) gives 2/7. This
  bound is weaker than the literature: Károlyi and Solymosi
  (arXiv:2607.22828, July 2026) cite an upper bound F(n) ≪ n²/(log
  n)\^{}(1+c) from Pilatte's extension of the Bloom--Sisask method,
  which implies the density tends to 0. They also give a lower bound
  F(n) = Ω(n\^{}1.3). Our block argument therefore adds nothing
  asymptotic.
\end{itemize}

\subsection{4. Limitations}\label{limitations}

\begin{enumerate}
\def\labelenumi{\arabic{enumi}.}
\tightlist
\item
  \textbf{Optimality is unknown for n ≥ 8.} The CP-SAT upper bounds (19,
  28, 40, 121, 144) are far above the best sets found. The gap for n =
  11 and 12 is large, so the lower bounds could be improved.
\item
  \textbf{Lower bounds for n ≥ 12 come from a single heuristic run per n
  (plus one CP-SAT run for n = 12).} Other seeds or longer runs may give
  larger values.
\item
  \textbf{Novelty check is limited.} We checked the OEIS entry for
  A271906 (which in fact contains lower bounds for n = 12--15) and one
  preprint (Károlyi--Solymosi, arXiv:2607.22828). The preprint's results
  concern asymptotics, not the small-n values in this paper. A published
  table for n ≥ 12 may exist, and we did not check further.
\item
  \textbf{No theoretical content.} We prove no general bound on a(n).
  Trivially, each 2×2 block of the grid holds at most two admissible
  points, since any three corners of an axis-aligned unit square form a
  right isosceles triangle. This gives a(n) ≤ 2⌈n/2⌉², which is far from
  the computed values.
\end{enumerate}

\subsection{5. Reproducibility}\label{reproducibility}

Files are in \texttt{./}:

\begin{itemize}
\tightlist
\item
  \texttt{rit.py}: triangle enumeration and CP-SAT model.
\item
  \texttt{local.py}: simulated annealing.
\item
  \texttt{verify.py}: independent checker.
\item
  \texttt{local\_n\{12..16\}.json}, \texttt{big\_n\{11,12\}.json}:
  reported sets.
\item
  \texttt{summary.json}: verification results.
\end{itemize}

Run \texttt{python3\ local.py\ N\ SECONDS\ SEED}, then
\texttt{python3\ verify.py}. The CP-SAT runs used 2 workers and time
limits of 900 s (n = 11) and 1500 s (n = 12). The machine had 2 cores.

\end{document}
